\documentclass[11pt,letterpaper]{article}
\usepackage[T1]{fontenc}
\usepackage[utf8]{inputenc}
\usepackage{lmodern,microtype}
\usepackage[margin=0.79in,headheight=15pt]{geometry}
\usepackage{parskip,enumitem,longtable,tabularx,booktabs,array,xcolor,fancyhdr,hyperref}
\definecolor{navy}{HTML}{17314C}
\definecolor{teal}{HTML}{087F78}
\hypersetup{colorlinks=true,linkcolor=navy,urlcolor=teal}
\setlength{\emergencystretch}{3em}
\setlist[itemize]{leftmargin=*,nosep,topsep=5pt}
\setlist[enumerate]{leftmargin=*,itemsep=3pt}
\renewcommand{\arraystretch}{1.22}
\pagestyle{fancy}\fancyhf{}
\fancyhead[L]{\small Enterprise Data Modeling}\fancyhead[R]{\small Documentation Suite}
\fancyfoot[C]{\thepage}
\newcommand{\srcbook}{\textit{The Data Model Resource Book}, revised edition, Volume 2}
\newcommand{\ResponseLine}{\par\vspace{1.05em}\noindent\rule{\linewidth}{0.4pt}\par}
\newcommand{\ResponseArea}{\ResponseLine\ResponseLine\ResponseLine}
\begin{document}
\begin{center}
{\LARGE\bfseries\color{navy} Manufacturing Digital Thread}\par
\vspace{0.25em}{\large Scenario-based enterprise case study}\par
\vspace{0.3em}{\small Original engineering study material \textbullet{} 2026 edition}
\end{center}
\vspace{0.8em}
\section{Case context}
A manufacturer makes a configurable product using separately versioned engineering parts, a BOM effective window, production runs and as-built serial records.


\section{Given business facts}
\begin{itemize}
\item Customer order creates a demand signal but not a completed asset.
\item Engineering revision 3 is released after a production order has been scheduled.
\item Approved substitute components may differ from the design BOM.
\item Defective serial numbers require tracing to production batches.
\end{itemize}

\section{Design and decision assignment}
\begin{enumerate}
\item Define product, part revision, BOM line, work order, serial inventory and consumption entities.
\item Explain effective revision selection and as-designed versus as-built query.
\item Propose a trace query from serial number to consumed lot and supplier.
\item Explain how rework and scrapping should avoid destructive historical edits.
\end{enumerate}

\section{Constraints and boundary risks}
\begin{itemize}
\item Every important state transition must carry origin, time and an accountable actor/service.
\item Changes must not silently erase the fact that was true when an agreement or event occurred.
\item Cross-system identifiers require mapping and collision/alias resolution; they are not presumed globally unique.
\item The architecture must describe ownership, sensitive information handling, failure recovery and at-least-once delivery.
\end{itemize}

\section{Deliverable matrix}
\begin{tabularx}{\linewidth}{@{}p{0.28\linewidth} p{0.68\linewidth}@{}}
\toprule
\textbf{Artifact} & \textbf{Acceptance criterion} \\ 
\midrule
Business vocabulary & Definitions approved by domain stakeholder; terms are unambiguous \\ 
Logical ERD & Cardinality, optionality and history rules consistent with scenarios \\ 
SQL schema & Constraints and transaction design map to narrative invariants \\ 
Architecture decision & Alternatives and consequences for one material trade-off \\ 
Quality evidence & At least six cases including three invalid or failure paths \\ 
\bottomrule
\end{tabularx}

\section{Facilitated review}
Ask the modeler to walk one event end-to-end and then introduce a late correction, partial failure and cross-system identifier conflict. Record any rule that the current design does not enforce; an unanswered question is an open architectural risk, not implicit approval.


\vfill
\hrule
{\footnotesize\textbf{Provenance.} Supplementary engineering practice: original content; the cited book provides domain-modeling context.}
\end{document}
